% smoke.tex — R1 smoke build driver, The Sourdough Framework: Arabic Edition
%
% Original: The Sourdough Framework by Hendrik Kleinwächter,
% the-sourdough-framework.com, CC BY-SA 4.0. The Arabic edition inherits
% CC BY-SA 4.0.
%
% Exercises the reconstructed pipeline end to end: cover + title page
% (cover-ar.tex, 4-word title), TOC with dotted numbers, chapter/section
% numbering LTR, bare Western digits in Arabic prose, \qty/\qtyrange in
% \LR{} islands, a mirrored table, a TikZ flowchart float, footnote with
% \textenglish parens, \pHvalue. Not a content unit: chapter text is
% re-typeset at R2 from the recovery references.

\documentclass[paper=a4, twoside=false, fontsize=12pt, parskip=half,
                bibliography=totoc, listof=totoc]{scrbook}
\usepackage{sourdough-ar}
\sisetup{range-phrase={\textendash}}
\author{Hendrik Kleinw\"achter}
\title{منظومة عمل الخبز الحامض}

\begin{document}
% cover-ar.tex — Arabic cover and title page for The Sourdough Framework
%
% Original: The Sourdough Framework by Hendrik Kleinwächter,
% the-sourdough-framework.com, CC BY-SA 4.0. The Arabic edition inherits
% CC BY-SA 4.0.
%
% Fork of the upstream cover/cover.tex. Page 1 carries the cover raster
% (upstream cover/cover-page.jpg, upstream commit 69636e6) with the Arabic
% overlay of record on the cover band; page 2 is the blank verso; page 3 is
% the typeset title page with the Arabic title, edition date and the
% author/translator credit.
%
% Overlay of record (Phase 26/27, Item 1 applied at R1: 4-word title):
%   Line 1: title, \put(88.45, 7.87), 14/16 Amiri bold, centered
%   Line 2: تأليف: هندريك كلاينفِختر · ترجمة: نائل العيله
%           \put(88.45, 2.92), 8/9 Amiri regular, centered
% The \makeatletter/\if@twoside wrapper is upstream-verbatim (a .tex file;
% the no-\makeatletter rule binds the .sty only).

\thispagestyle{empty}
\setlength{\unitlength}{1mm}
\makeatletter
\if@twoside%
    \noindent\begin{picture}(0,0)(1,-1)
    \put(-16.3,-265){\includegraphics[width=1.33\linewidth]{cover/cover-page.jpg}}
    \end{picture}
\else%
    \noindent\begin{picture}(0,0)(1,-1)
    \put(-23.3,-265){\includegraphics[width=1.33\linewidth]{cover/cover-page.jpg}}
    \end{picture}
\fi%
% Arabic overlay on the cover band (Phase 26/27 of record; 4-word title).
\noindent\begin{picture}(0,0)(1,-1)
    \put(88.45,7.87){\makebox(0,0)[c]{\fontsize{14}{16}\selectfont\bfseries منظومة عمل الخبز الحامض}}%
    \put(88.45,2.92){\makebox(0,0)[c]{\fontsize{8}{9}\selectfont تأليف: هندريك كلاينفِختر · ترجمة: نائل العيله}}%
\end{picture}
\makeatother

\newpage
\thispagestyle{empty}

\rule{1pt}{\textheight} % Vertical line
 % Whitespace between the vertical line and title page text
\hspace{0.05\textwidth}
 % Paragraph box for holding the title page text, adjust the width to move the
% title page left or right on the page
\parbox[b]{0.75\textwidth}{%
{\Huge\bfseries منظومة عمل الخبز الحامض}\\[1\baselineskip] % Title (4-word form, Item 1 of record)
{\large\textenglish{The Sourdough Framework}}\\[2\baselineskip] % Original-title subtitle, accepted p. 2
{\large\textit{الإصدار: 7 نوفمبر 2026}}\\[4\baselineskip] % date frozen to the accepted edition's build date
{\Large هندريك كلاينفِختر \textenglish{(Hendrik Kleinwächter)}}\\[1\baselineskip]
{\large ترجمة: نائل العيله}

% Whitespace between the title block and the copyright text. In the PDF and
% the ebook it holds a note on asking this book: Loafy answers from it. The
% website leaves it out, it has Loafy on every page already. Each place has
% its own via tag, so Loafy can count where readers came from.
\ifdefined\HCode
\vspace{0.5\textheight}
\ifdefined\sdfebook
\textbf{أسئلة أثناء الخَبز؟} اسأل هذا الكتاب. \textenglish{Loafy}، رفيقه المجاني،
يجيب من هذه الصفحات ويسمّي القسم:
\href{https://loafy.ai/loafy?via=tsf-epub-title}{https://loafy.ai}
\fi
\else
\vspace{0.17\textheight}

\noindent
\begin{minipage}[c]{2.4cm}
% Black, not link blue: a light QR code is hard to scan once printed.
% The QR matrix macros must typeset LTR; RTL context corrupts them.
\LRE{\qrcode[nolink,height=2.2cm]{https://loafy.ai/loafy?via=tsf-pdf-qr}}
\end{minipage}%
\hspace{0.4cm}%
\begin{minipage}[c]{\linewidth-2.8cm}
\textbf{أسئلة أثناء الخَبز؟} اسأل هذا الكتاب «لُوفي» \textenglish{(Loafy)}، رفيقه المجاني،
يجيب عنك من صفحاته ويسمّي لك القسم:
\href{https://loafy.ai/loafy?via=tsf-pdf-title}{\nolinkurl{https://loafy.ai}}
\end{minipage}

\vspace{0.17\textheight}
\fi


{\noindent
\begin{flushleft}
    \includegraphics[width=3cm]{cover/CC-BY-SA}\par
الشيفرة المصدرية الكاملة لهذا الكتاب متاحة على
\url{https://github.com/hendricius/the-sourdough-framework/} برخصة
\textenglish{CC-BY-SA}.
راجع \url{https://creativecommons.org/licenses/by-sa/4.0/} لمزيد من التفاصيل.
لا تتردد في الإبلاغ عن الأخطاء أو اقتراح التحسينات. يتوفر إصدار غلاف مقوّى من
الكتاب أيضًا. مزيد من المعلومات هنا:
\url{https://www.breadco.de/hardcover-book}

الإصدار العربي: مبنيّ على الكتاب الأصلي \textenglish{The Sourdough Framework}
لـ هندريك كلاينفِختر، المنشور مرخّصًا برخصة \textenglish{CC BY-SA 4.0}،
ويخضع هذا الإصدار العربي للرخصة نفسها.
\end{flushleft}
}

{\centering\textenglish{Original: The Sourdough Framework by Hendrik
Kleinwächter, the-sourdough-framework.com, CC BY-SA 4.0}\par}
}


\frontmatter
{\hypersetup{hidelinks}\tableofcontents}

\mainmatter
\setchapterstyle{lines}

\chapter{فصل اختبار البنية التحتية}

هذا فصل اختبار للبنية التحتية للطبعة العربية. يجرّب الأرقام الغربية
في السياق العربي: تُخمَّر العجينة مدة 24 ساعة عند درجة حرارة 20 درجة،
ثم يُخبز الرغيف مدة 45 دقيقة. النسبة المئوية للترطيب هي 75% من وزن
الدقيق، ودرجة الحرارة المثالية 26.5 درجة. راجع الفصل 3 والمخطط 1.2
للتفاصيل.

\section{القسم الأول: الأرقام والوحدات}

تحتفظ خلايا \qty{500}{\gram} من الدقيق بـ\qty{350}{\milli\liter} من
الماء\footnote{الأقواس اللاتينية تمر عبر \textenglish{(parenthetical
English)} إلى الخط اللاتيني الرئيسي.}، ويعمل الخليط عند
\qtyrange{20}{25}{\celsius}. قيمة الرقم الهيدروجيني تهبط إلى
\pHvalue{3.8} بعد التخمير الكتلي. نسبة 1:2 تعني جزءًا من البادئ
وجزءين من الدقيق، والمدى 10--15% يغطي الفروق الموسمية.

\subsection{قسم فرعي مرقّم}

الخلية الواحدة تحتوي نصًا عربيًا يتدفق من اليمين إلى اليسار، بينما
يبقى الرقم 42 مرتبًا من اليسار إلى اليمين داخل السطر.

\begin{table}[!htb]
\centering
\caption{جدول اختبار المرآة}
\label{tab:smoke-mirror}
\begin{tabular}{lll}
\toprule
\thead{\qty{250}{\gram}} & \thead{الوصف} & \thead{العنصر} \\
\midrule
\qty{500}{\gram} & دقيق قمحي كامل الحبة & الدقيق \\
\qty{350}{\milli\liter} & ماء بدرجة حرارة الغرفة & الماء \\
\qty{100}{\gram} & بادئ حامض نشط & البادئ \\
\qty{10}{\gram} & ملح خشن & الملح \\
\bottomrule
\end{tabular}
\end{table}

\begin{flowchart}[!htb]
\centering
\begin{tikzpicture}[node distance = 3.5cm, auto]
  \node [start] (start_n) {\textarabic{اخلط \qty{50}{\gram} دقيق مع \qty{50}{\gram} ماء}};
  \node [block, right of=start_n, node distance = 3.5cm] (wait_n) {\textarabic{انتظر\\ \qty{24}{\hour}}};
  \node [decision, right of=wait_n, node distance = 3.5cm] (readycheck_n) {\textarabic{جاهز؟}};
  \draw [line] (start_n) -- (wait_n);
  \draw [line] (wait_n) -- (readycheck_n);
  \draw [line] (readycheck_n) -- node {\textarabic{نعم}} ++(3.5,0);
\end{tikzpicture}
\caption{مخطط انسيابي اختباري}
\label{fig:smoke-flowchart}
\end{flowchart}

\end{document}
